\section{Тестирование}
Вопросы к системе представлены в листингах~\ref{lst:test1}--\ref{lst:test12}.

\begin{listing}[Поиск всех бабушек субъекта]{lst:test1}{Prolog}
?- find_grandparent(GP, female, _, svetlana).

GP = maria ;
GP = olga .
\end{listing}

\begin{listing}[Поиск всех дедушек субъекта]{lst:test2}{Prolog}
?- find_grandparent(GP, male, _, svetlana).

GP = ivan ;
GP = petr .
\end{listing}

\begin{listing}[Поиск всех бабушек и дедушек субъекта]{lst:test3}{Prolog}
?- find_grandparent(GP, _, _, svetlana).

GP = ivan ;
GP = maria ;
GP = petr ;
GP = olga .
\end{listing}

\begin{listing}[Поиск бабушки по материнской линии]{lst:test4}{Prolog}
?- find_grandparent(GP, female, female, svetlana).

GP = olga .
\end{listing}

\begin{listing}[Поиск бабушки и дедушки по материнской линии]{lst:test5}{Prolog}
?- find_grandparent(GP, _, female, svetlana).

GP = petr ;
GP = olga .
\end{listing}

\begin{listing}[Вычисление максимума из двух чисел с использованием отсечения]{lst:test6}{Prolog}
?- max2_cut(15, 42, Max).

Max = 42.
\end{listing}

\begin{listing}[Вычисление максимума из двух чисел без использования отсечения]{lst:test8}{Prolog}
?- max2(15, 42, Max).

Max = 42.
\end{listing}

\begin{listing}[Вычисление максимума из двух чисел с использованием отсечения]{lst:test9}{Prolog}
?- max2_cut(42, 42, Max).

Max = 42.
\end{listing}

\begin{listing}[Вычисление максимума из трех чисел без использования отсечения]{lst:test10}{Prolog}
?- max3(7, 24, 13, Max).

Max = 24.
\end{listing}

\begin{listing}[Вычисление максимума из трех чисел с использованием отсечения]{lst:test11}{Prolog}
?- max3_cut(24, 7, 13, Max).

Max = 24.
\end{listing}

\begin{listing}[Вычисление максимума из трех чисел с использованием отсечения]{lst:test12}{Prolog}
?- max3_cut(12, 12, 12, Max).

Max = 12.
\end{listing}